\documentclass[../ppgc-tese.tex]{subfiles}
\begin{document}

\begin{flushright}
\mbox{}%\vfill
{\sffamily\itshape
``Not all those who wander are lost.'' \\}
--- \textsc{J.R.R.~Tolkien}
\end{flushright}

This chapter studies one learning phenomenon, late generalization or grokking, and shows that it can be induced on demand by a controlled distribution shift between training and test data. The evidence is behavioral: it rests on training and test curves over synthetic datasets whose geometry we control, not on measurements of the models' internal representations.

\section{Grokking: Late Generalization}
\label{sec:grokking}
The grokking phenomenon is a sudden improvement in model performance on unseen data after a period of seemingly stagnant learning, also known as late or delayed generalization. It was first detected in 2021, in the work of \cite{Grokking}. The subsequent literature attributed the phenomenon mostly to three main factors: data sparsity, large initialization values, and high regularization rates. Although the combination of these factors may indeed cause grokking in some situations, they cannot be said to be the cause or to provide a complete explanation of grokking.

As illustrated in Figure \ref{fig:grokking-def}, grokking can be characterized by the transitioning from a low-performing region to a high-performing region as measured on the test set (curve $T_E$ in Fig.\ref{fig:grokking-def}) after the training set performance has converged ($T_R$ in Fig.\ref{fig:grokking-def}). The interval between the inflection points $\alpha_0$ (end of memorization phase) and $\alpha_1$ (start of generalization phase) signifies the transition phase, during which the test set performance gradually improves and aligns with the performance on the training set. 

Understanding intriguing learning phenomena such as grokking can provide valuable insights into the inner workings of deep neural networks. Late generalization challenges traditional stopping criteria, which rely on monitoring training and test set losses. While it remains difficult to predict when late generalization might occur, this chapter advances the discussion by elucidating how distribution shifts contribute to grokking. Specifically, we explore how modeling class relationships can increase the probability of generalization and guide the design of alternative stopping criteria. These findings help identify conditions under which models are more likely to generalize after training convergence.

\begin{definition}
[Grokking]\label{grokking-definition}    
Let $\Delta_T$ denote an interval with lower bound $\sigma_E$ containing the limit value of a converging training set performance curve at time $\alpha_0$. Let $\sigma_S \ll \sigma_E$ denote an upper bound on the test set performance measured up to time point $\alpha_0$. We say that \textit{grokking} took place if at time point $\alpha_1$ > $\alpha_0$, the test set performance is within $\Delta_T$ and remains within $\Delta_T$ for an interval $S \gg (\alpha_1 - \alpha_0)$.
\end{definition}

\begin{figure}[ht]
    \centering
    \includegraphics[width=0.7\linewidth]{figures/grokking_definition.pdf}
    \caption{An illustration of the key terms used to define grokking. After training accuracy ($T_R$) converges, test accuracy ($T_E$) transitions from a low-performance plateau (ending at $\alpha_0$) to a high-performance one (starting at $\alpha_1$). We measure the grokking intensity by the performance gain, $\Delta_S$, and its stability by the support interval, $S$.}
    \label{fig:grokking-def}
\end{figure}

\subsection{Grokking: Visualization and metrics}

Our study reveals that a specific type of distribution shift, caused by the sampling of data organized into classes and subclasses, can reliably produce late generalization. Unlike prior research, where grokking was observed sporadically, we demonstrate that grokking occurs systematically when certain subcategories or subpatterns in the data are under-represented. Our experiments suggest that a model's ability to generalize to unseen data depends on subtle data interactions, which may emerge under weak training signals, leading to significant improvements in test performance --- even when performance on specific subcategories remains poor.

The contributions of this chapter are as follows:
\begin{enumerate*}[label=(\roman*)]
    \item We introduce two synthetic datasets based on class hierarchies, designed to serve as benchmarks for analyzing late generalization and grokking intensity; 
    \item We replicate previous experiments and conduct new analyses, showing that grokking is more strongly associated with distribution shifts than with sparsity and regularization, as previously believed \footnote{The cripts are provided for replication in \url{\repourl}.};
    %\item We conduct experiments on the MNIST dataset, where we induce distribution shifts via clustering of latent representations, examining the limits of our findings and providing a case study on the challenges of reproducing the phenomenon on complex real-world data;
    \item We report results using both multilayer perceptrons and transformers, concluding that late generalization is largely architecture-independent. 
\end{enumerate*}


However, if the evaluation loss decreases sharply at a later time, even after the training loss has approached zero, this indicates a capability to generalize, rather than memorization, given a very weak training signal that needs to be understood. 

\subsection{Distribution-shift impact on equidistant subclasses}
\begin{figure*}[!ht]
    \centering
    \begin{subfigure}[b]{\textwidth}
         \centering
         %\includegraphics[width=1.06cm]{figures/legend.pdf}
         
         \includegraphics[width=\textwidth]{figures/equidistant/result_eq_line-1.pdf}
         \caption{Effect of dataset size on the equidistant dataset. Grokking occurs across all sample sizes, but its intensity ($\Delta_s$) is inversely proportional to the number of samples.}
         \label{fig:balanced_data_varying_sample_size}
    \end{subfigure}%\vspace{0.5cm}

    \begin{subfigure}[t]{\linewidth}
         \centering
         \includegraphics[width=0.6\textwidth]{figures/equidistant/result_small.pdf}
         \caption{Effect of subclass removal on the equidistant dataset. Removing a subclass ($f=0$, top) prevents generalization, while re-introducing a small fraction ($f=0.2$, bottom) is sufficient to induce grokking.}
         \label{fig:small_unbalance_balance_eq}
    \end{subfigure}
    \hfill
    \begin{subfigure}[t]{\linewidth}
         \centering
         \includegraphics[width=0.6\textwidth]{figures/equidistant/result_large.pdf}
         \caption{Grokking in large datasets. With a highly skewed distribution ($f \ll 1$), the phenomenon is observed even in larger datasets, contrary to what was previously assumed.}
         \label{fig:medium_unbalance_balance}
    \end{subfigure}

    \caption{Grokking dynamics in the equidistant dataset under different sampling conditions. (a) The effect of varying dataset size with balanced sampling. (b) Comparing subclass removal ($f=0$) with under-sampling ($f=0.2$). (c) Inducing grokking in large datasets with extreme imbalance.}
    \label{fig:exp_samples_100-200_frac_0.0-0.2-1.0}
\end{figure*}


%This phenomenon challenges conventional wisdom and underscores the importance of evaluating model performance not just on the training data but also on unseen data to ensure robust generalization capabilities.

Our hypothesis is that grokking is a statistical phenomenon derived from a distribution shift between test data and training data. This difference allows for the regularization to shift the network from purely a local minimum of the test loss surface to a minimum of the test loss as well which can be reproduced systematically.

\textbf{Distribution shifts inducing grokking:}
% \begin{itemize}
%     \item Loss function -> strong incentive to keep decision boundaries
%     \item Regularization as a tool to reach a simpler minimum (less parameters)
% \cite{Grokking} showed grokking in models that do not rely on an internal latent space to represent the data points which points that although the particular model is relevant the data distribution is a more determinant factor.
% \end{itemize}
The phenomenon of grokking, induced by distribution shifts, can be attributed to several key factors.
Firstly, even a faint, i.e. close to zero, training signal may interact with related data within a class or across nearby classes to provide an incentive for models to maintain decision boundaries.
Additionally, regularization serves as a tool to guide models toward simpler solutions with a sparse weight matrix influencing but not determining the occurrence of grokking.
Research by \cite{Grokking} demonstrated grokking in models that do not rely on an internal latent space to represent data, suggesting that while the specific model architecture may be relevant, it is the underlying data distribution and the compounding of learning that will determine late generalization.


In order to confirm or reject our hypothesis, we created two synthetic datasets that allow for controlling the assumed underlying conditions for grokking, as described next. 

%We then... NIST data set... to test those conditions in a more realistic setting.


\section{New Synthetic Datasets for Grokking}\label{sec:artificial-datasets}

The two synthetic datasets introduced in this chapter enable the induction of grokking with minimal tuning of model parameters, avoiding an extensive search of the parameter space. They allow us to create a controlled shift between the training and test distributions and to decide how much information each subclass carries, as detailed below. Both datasets are released with the code accompanying this thesis\footnote{\url{\repourl}}, and Appendix~\ref{app:dataset-creation} documents their construction step by step.

The primary goal of the datasets is to be linearly separable when considering all sample dimensions $d$, while also being non-trivial to solve under noise or extreme down-sampling. The datasets are defined by projecting classes onto a $d$-dimensional space using multivariate normal distributions, ensuring control over both class definitions and their distances in that space.

We define classes $C_1, \dots, C_n$, where each class $C_i$, $1 \leq i \leq n$, is the union of the samples in the disjoint subclasses $C_{i,1}, \dots, C_{i,m}$. We say that $C_i = \bigcup_{j=1}^{m} C_{i,j}$, and $C_{i,p} \cap C_{i,k} = \emptyset$ for all $k \neq p$. The classification target is always the class $C_i$; subclasses are never shown to the model and exist only to define the geometry of the data and to control the sampling. The above structure of classes and subclasses mimics the organization of some real-world datasets, such as COCO \cite{lin2014COCO}, which is used in image segmentation where objects can be part of a larger scene or other objects. More importantly, it provides a straightforward and effective method to induce distribution shifts on demand by down-sampling specific subclasses while preserving the original proportions of the other classes.

\noindent\textbf{Design rationale.}
Our hypothesis (Section~\ref{sec:grokking}) states that a faint training signal may interact with related data \emph{within a class or across nearby classes} to sustain a decision boundary that the loss alone no longer rewards. This statement has two separable ingredients: (i) the existence of a small number of samples from the shifted region of the input space, which provides the faint signal itself; and (ii) the existence of retained training data that is \emph{related} to that region, which provides something for the faint signal to interact with. The two datasets are designed to isolate ingredient (ii) while keeping everything else fixed. They share the same generator, the same dimensionality $d = 117$, the same noise model, the same four-way classification target, and the same sub-sampling protocol; the only manipulated variable is the geometry of the subclasses relative to the classes. In the first dataset, every subclass is equidistant from every other subclass, so removing a subclass leaves its region of the input space without any preferential neighbor in the training set: the region can only be recovered from the samples of the removed subclass itself. In the second dataset, the subclasses of a class are strictly closer to each other than to the subclasses of any other class, so the retained subclasses of a class always \emph{cover} the region of a removed sibling. We refer to this second situation as \emph{proximity-mediated information transfer}, and define it operationally in Table~\ref{tab:removed-neighbors}: given a removed subclass, we compare the distance from its centroid to the nearest retained centroid of the same class with the distance to the nearest retained centroid of any other class. Transfer is possible when the former is strictly smaller; it is absent when the two are equal.

\noindent\textbf{Common construction.}
Both datasets are produced by the same generator, which we summarize here and detail in Appendix~\ref{app:dataset-creation}. Let $p = 9$ be the number of \emph{dimension groups} and $r = 13$ the number of dimensions per group, so that $d = pr = 117$. A subclass is identified with a binary code $c \in \{0,1\}^p$, and we write $b_\ell(c)$ for its $\ell$-th bit. For every group $\ell$ we draw, once, two centers $\mu_\ell^0 \sim \mathcal{U}[-1,1]^r$ and $\mu_\ell^1 = \mu_\ell^0 + u_\ell$ with $\|u_\ell\| = 1$, and a single covariance $\Sigma_\ell = R_\ell^\top S_\ell^2 R_\ell$, where $S_\ell$ is diagonal with entries $0.125\,(1 \pm 0.25\,\xi)$, $\xi \sim \mathcal{U}[0,1]$, and $R_\ell$ is a product of random Givens rotations with angles in $\pm[\pi/6, \pi/3]$. A sample of subclass $c$ is the concatenation $x = (x_1, \dots, x_p)$ of $p$ independent blocks,
\begin{equation}\label{eq:generator}
    x_\ell = \mu_\ell^{\,b_\ell(c)} + \eta_\ell S_\ell R_\ell,
    \qquad \eta_\ell \sim \mathcal{N}(0, I_r),
    \qquad \ell = 1, \dots, p.
\end{equation}
Because the two centers of a group share the same covariance and $\|\mu_\ell^1 - \mu_\ell^0\| = 1$ for every group, the Euclidean distance between the centroids of two subclasses is determined by the Hamming distance $H(\cdot,\cdot)$ between their codes:
\begin{equation}\label{eq:centroid-distance}
    \|\centroid{C_a} - \centroid{C_b}\| = \sqrt{H(a, b)}.
\end{equation}
This is the property that lets us specify the geometry of a dataset by choosing codes. Choosing a set of codes at constant mutual Hamming distance yields equidistant subclasses; choosing a code together with its Hamming-1 neighbors yields a tight cluster of subclasses. The within-subclass noise is small relative to the unit separation between centers (per-axis standard deviations range between $0.10$ and $0.16$ in the released files), so each pair of subclasses is separable by a hyperplane with negligible overlap, and so are the classes defined below. The task is therefore easy when the training set is representative, but sensitive to which subclasses are present in it.

\noindent\textbf{Sub-sampling protocol.}
We cause arbitrary distribution shifts in the datasets by selecting a specific set of subclasses to contain fewer samples than the other subclasses. Without loss of generality, we create all subclasses with the same number of samples in every dataset that we study, i.e. $|C_{i,j}| = s$, for all $i,j$. To down-sample a subclass we define $f \in [0, 1]$ such that if $f = 0$ then every sample of that subclass is removed and if $f = 1$ no change is made to its proportion. Let $\gamma_D$ denote the total number of training samples to draw, $\gamma_s$ the number of subclasses to sub-sample, and $\gamma_r$ the number of remaining subclasses. Each sub-sampled subclass receives weight $f$ and each remaining subclass receives weight $1$, and the number of samples drawn from each is proportional to its weight:
\begin{equation}\label{eq:num_samples_unbalanced}
    s_{s} = \operatorname{round}\!\left(\frac{f\gamma_D}{f\gamma_s + \gamma_r}\right)
    \quad \text{and} \quad
    s_{r} = \operatorname{round}\!\left(\frac{\gamma_D}{f\gamma_s + \gamma_r}\right).
\end{equation}
Whenever rounding leaves a residual $\gamma_D - (\gamma_s s_s + \gamma_r s_r) \neq 0$, it is distributed one sample at a time to the subclasses with the largest (or smallest) rounding remainders, so that exactly $\gamma_D$ samples are drawn. For example, with four classes of two subclasses each, $\gamma_D = 2000$, $f = 0.2$, and one sub-sampled subclass per class ($\gamma_s = \gamma_r = 4$), we obtain $s_{s} = \operatorname{round}(0.2 \cdot 2000 / 4.8) = 83$ and $s_{r} = \operatorname{round}(2000 / 4.8) = 417$ samples, which sum to exactly $2000$.

In every experiment reported in this chapter, the sub-sampled subclasses are half of the subclasses of \emph{each} class, chosen once at random with a fixed seed and kept fixed across runs, so that $\gamma_s = \gamma_r$. Training sets of size $\gamma_D$ are drawn from a held-out training pool according to Equation~\ref{eq:num_samples_unbalanced}, whereas the test set is a balanced sample of $10{,}000$ examples in which every subclass is equally represented. The distribution shift is therefore entirely on the training side, and $f$ controls its severity: $f = 1$ is the balanced case and $f = 0$ removes the sub-sampled subclasses altogether. We study equidistant subclasses (Dataset~1) and equivariant subclasses (Dataset~2), as described next.

\begin{figure*}[t]
    \centering

    \begin{subfigure}[t]{0.475\textwidth}
         \centering
         \includegraphics[width=0.3\textwidth]{figures/group_eqdist.pdf}
         \caption{A schematic of the equidistant dataset. Each subclass is equidistant from all others, including its sibling in the same class, so a removed subclass has no preferential neighbor in the training set.}
         \label{fig:group_eqdist}
    \end{subfigure}
    \hfill
    \begin{subfigure}[t]{0.475\textwidth}
         \centering
         \includegraphics[width=0.4\textwidth]{figures/group_eqvariant.png}
         \caption{A schematic of the equivariant dataset. Subclasses are grouped by proximity into larger classes, so the retained subclasses of a class always cover the region of a removed sibling.}
         \label{fig:group_eqvariant}
    \end{subfigure}
    \caption{Schematic representations of the (a) equidistant and (b) equivariant synthetic datasets developed for this study.}
    \label{fig:synthetic_dataset}
\end{figure*}

\noindent \textbf{Dataset 1 (equidistant subclasses)}\label{sec:eqdist-dataset}: this dataset is the control regime, in which no proximity-mediated transfer is possible. It has $4$ classes, $C_1,\dots, C_4$, each composed of two disjoint subclasses, $C_i = C_{i,1} \cup C_{i,2}$. The eight subclasses are the codes of a maximum clique of the graph that connects codes at Hamming distance exactly $2$, so by Equation~\ref{eq:centroid-distance} the centroids of all subclasses --- no matter which class they belong to --- are at the same distance $\sqrt{2}$ from one another. The classes are formed by pairing these codes, and the resulting class centroids are also equidistant, at distance $1$ from one another. This concept is illustrated in Figure~\ref{fig:group_eqdist}. As a consequence, when one subclass per class is removed from the training set, the nearest retained centroid is at distance $\sqrt{2}$ both within the same class and in every other class (Table~\ref{tab:removed-neighbors}). Nothing in the retained data indicates to which class the removed region belongs, and its assignment is left entirely to the inductive bias of the model and the optimizer. This dataset therefore tests whether late generalization requires at least a faint signal from the shifted region itself: if it does, we expect $f = 0$ to leave the model on a plateau and small $f > 0$ to produce grokking. The released files contain $2{,}500$ samples per subclass ($20{,}000$ per split).

\noindent \textbf{Dataset 2 (equivariant subclasses)}\label{sec:similar-dataset}: this dataset is the transfer regime. It has $4$ classes, each composed of $10$ subclasses. The four classes are built around four \emph{seed} codes at mutual Hamming distance $6$; each class consists of its seed together with the nine codes at Hamming distance $1$ from it. By Equation~\ref{eq:centroid-distance}, within a class the seed subclass is at distance $1$ from each of the other nine, which are at distance $\sqrt{2}$ from one another; across classes, subclass centroids are at distance at least $2$ (Hamming distance at least $4$). That is, if $\centroid{C_{i,j}}$ is the centroid of subclass $j$ of class $i$, then
\[
\|\centroid{C_{i,j}} - \centroid{C_{i,l}}\| \leq \sqrt{2} < 2 \leq \|\centroid{C_{i,j}} - \centroid{C_{k,l}}\|
, \quad i \neq k,
\]
for all $j, l$. Subclasses within a class are thus closer to each other than to any subclass of another class, but, unlike in Dataset~1, they are not all equidistant: each class has a hub-and-spoke structure, illustrated in Figure~\ref{fig:group_eqvariant}. When half of the subclasses of a class are removed, the nearest retained centroid is always in the same class --- at distance $1$ if the seed is retained (or if the removed subclass is the seed itself) and $\sqrt{2}$ otherwise --- while the nearest retained centroid of any other class is at distance at least $2$ (Table~\ref{tab:removed-neighbors}). In the sub-sampling used in our experiments the seed is among the removed subclasses in two of the four classes, so both strengths of coverage occur. This dataset therefore tests whether the retained relational structure can substitute for the missing samples: if it can, we expect generalization, and possibly late generalization, even at $f = 0$. The released files contain $500$ samples per subclass ($20{,}000$ per split).

\begin{table}[t]
    \centering
    \caption{Distance from the centroid of each removed subclass to the nearest retained centroid of the same class and of any other class, for the sub-sampling used in the experiments (half of the subclasses of every class). ``Removed'' lists which subclasses of the class are removed; ``others'' are non-seed subclasses. Values follow from Equation~\ref{eq:centroid-distance} and match the released files to two decimal places. Proximity-mediated transfer is possible when the same-class distance is strictly smaller.}
    \label{tab:removed-neighbors}
    \small
    \begin{tabular}{llcccl}
        \toprule
        Dataset & Class & Removed & Same class & Other class & Transfer \\
        \midrule
        Equidistant & each of $C_1$--$C_4$ & 1 of 2 & $\sqrt{2}$ & $\sqrt{2}$ & no \\
        \midrule
        \multirow{3}{*}{Equivariant} & \multirow{2}{*}{$C_1$, $C_2$ (seed removed)} & seed & $1$ & $\sqrt{5}$ & yes \\
         & & 4 others & $\sqrt{2}$ & $2$ & yes \\
         & $C_3$, $C_4$ (seed retained) & 5 others & $1$ & $2$ & yes \\
        \bottomrule
    \end{tabular}
\end{table}

\noindent\textbf{Limitations of the comparison.}
The two datasets differ not only in geometry but also in the number of subclasses per class ($2$ versus $10$) and, consequently, in the number of retained subclasses per class after sub-sampling ($1$ versus $5$). A model trained on Dataset~2 therefore sees a more diverse support for each class, and part of the difference in behavior at $f = 0$ could in principle be attributed to this diversity rather than to proximity alone. The geometry of Dataset~1 rules out proximity-mediated transfer regardless of how many subclasses are retained, so the control is sound; disentangling diversity from proximity in the transfer regime would require an equivariant variant with two subclasses per class, which we leave for future work.

\section{Experimental Results}\label{sec:results}

All our experiments were designed to run on a single GPU, although we ran multiple experiments in parallel to expedite the data analysis process. Each compute node used in the experiments was equipped with an NVIDIA V100 GPU and 16 GB of memory. We conducted experiments in batches, varying the sample size, fraction $f$, and model parameters such as the learning rate and weight decay. Each batch took less than 24 hours to run under these settings. For robustness sake, the results presented in this chapter are aggregated from 10 independent runs. 

\subsection{Network architectures evaluated}\label{sec:archtectures}
 % \begin{itemize}
 %     \item MLP
 %     \item 2-layer Transformer
 %     \item Throughout the paper we are showing results for the MLP architecture since the results on both architectures showed very similar behavior. For completeness, the reader might refer to the appendices~\ref{app:equivariant_dataset}, and~\ref{app:equidistant_dataset} for the results on the 2-layer transformer architecture.
 %     \item The start of each plot in this paper occurs after 100 optimization steps, this is usually enough time for the algorithm converge to perfect train accuracy.
 %     \item We use a batch size of 200 samples in all our experiments, except in the cases we downsample a training dataset to 100 samples to investigate its impact in the grokking phenomenon, in this case we use batch size of 100 samples.
 % \end{itemize}

%We evaluated different architectures on each dataset to inspect how training behaves and under which conditions grokking occurs.

We conducted our experiments using two different architectures: a Multi-layer Perceptron (MLP) and a 2-layer Transformer.
Throughout this chapter, we primarily present results for the MLP architecture since the results obtained with both architectures exhibited very similar patterns.
For completeness, the results obtained with the 2-layer Transformer can be found in the git repository. %Appendix~\ref{app:dataset-creation}.
We also provide the code for replicating the experiments in the repository \url{\repourl}.

% Each plot in this paper starts after 100 optimization steps, which is typically sufficient for the algorithm to converge to perfect training accuracy. In all our experiments, we use a batch size of 200 samples.
%
In the experiments reported in this chapter we set the learning rate of the models to $10^{-3}$, the weight decay to $10^{-4}$, and have an initialization scale factor of 8 (i.e., we multiply all the weights of the model to this number as described in \cite{liu2022Grokking}, for both architectures tested.

\textbf{Distribution-shift caused by data sparsity} When performing a balanced down-sampling of our training set, i.e. without forcing any subclass to be sampled more than any other, we find that grokking diminishes as the number of samples increases, only existing in considerably small sample sizes, as depicted in Figure~\ref{fig:balanced_data_varying_sample_size}. This is consistent with the literature, as discussed in Section~\ref{sec:grokking}.

\textbf{Distribution-shift caused by unbalanced equidistant subclass sampling}
In order to evaluate the impact of distribution shifts in the  occurrence of grokking, we evaluated different data unbalancing created by the use of different values for $f$ as described in Section~\ref{sec:artificial-datasets}. It is clear from Figure~\ref{fig:small_unbalance_balance_eq} that removing completely a subclass has a detrimental effect on generalization, with the models being stuck on a sub-optimal plateau. Nevertheless, adding even a very small number of examples in that subclass (as defined by the fraction e.g. $f = 0.01$ of the missing subclasses) not only allowed the model to generalize well, but allowed the model to generalize late (i.e. grokking occurred). Since this very small fraction of the examples was not visibly relevant for training (was not misclassified) their impact on the test set performance was seen only after training converged and weight decay led the network to a more sparse internal representation.

When we increase the training data size, as mentioned, we see a diminishing effect on grokking. In order to keep analyzing this, we then evaluated more aggressive imbalance factors, such as $f = 0.01$ and $f = 0.05$. In a nutshell, Figure~\ref{fig:medium_unbalance_balance} shows that the data distribution shift alone is able to induce late generalization.

% -----------------
\subsection{Distribution-shift impact on equivariant subclasses}
We investigate the effects of sampling strategies on the grokking phenomenon by examining unbalanced sampling of classes. Firstly, we analyze balanced sampling of classes, as illustrated in Figure~\ref{fig:shift-equifariant}. Next, we explore the impact of unbalanced sampling, depicted in Figure~\ref{fig:equivariant-few-samples}. Our findings reveal proximity-mediated information transfer among the subclasses of a class, in the sense defined in Section~\ref{sec:artificial-datasets} and Table~\ref{tab:removed-neighbors}, which is evidenced by the comparison between the similar graphs with and without sampling (Figures~\ref{fig:equivariant-no-sample} and \ref{fig:equivariant-few-samples}).

\textbf{Distribution-shift caused by unbalanced equivariant subclass sampling}
The other setting evaluated in our experiments consisted of inducing different imbalances, using different values of $f$ as described in Section~\ref{sec:artificial-datasets}. The first noteworthy result here is that we were able to detect grokking even in the cases where samples were removed completely from specific subclasses, as depicted in Figure~\ref{fig:equivariant-few-samples}. This surprising behavior contrasts with the results obtained for the equidistant dataset. It is attributed to the fact that in the equivariant case, subclasses of a given class are closer to each other than to subclasses of other classes (Table~\ref{tab:removed-neighbors}): the retained subclasses of a class cover the region of a removed sibling and thus carry information that may result in grokking, whereas in the equidistant case the removed region is equally close to every class.

\begin{figure*}[t]
    \centering
    \begin{subfigure}[b]{\linewidth}
         \centering
         %\includegraphics[width=1.06cm]{figures/legend.pdf}
         \includegraphics[width=0.67\textwidth]{figures/result_exp_unb_samples_100_400_800_frac_0.0-1.0a.pdf}
         \caption{Generalization occurs even if a subclass is removed ($f=0$).}
         \label{fig:equivariant-no-sample}
    \end{subfigure}
    
    \vspace{0.5cm}
    
    \begin{subfigure}[b]{\linewidth}
         \centering
         \includegraphics[width=5.5cm]{figures/legend.pdf}
         
         \includegraphics[width=0.67\textwidth]{figures/result_exp_unb_samples_100_400_800_frac_0.0-1.0b.pdf}
         \caption{An imbalanced sample ($f=0.2$) induces a more classic grokking pattern.}
         \label{fig:equivariant-few-samples}
    \end{subfigure}
    \caption{Grokking in the equivariant dataset, where subclasses are structurally similar. Unlike in the equidistant dataset, the model generalizes even when a subclass is absent ($f=0$, top plots).}
    \label{fig:shift-equifariant}
\end{figure*}


% \section{Double descent}

\end{document}



